\documentclass[11pt]{article}
\usepackage[margin=1in]{geometry}
\usepackage[T1]{fontenc}
\usepackage{lmodern,amsmath,amssymb,amsthm,mathtools,booktabs,longtable,microtype}
\usepackage[hidelinks]{hyperref}
\setlength{\emergencystretch}{3em}
\allowdisplaybreaks
\newtheorem{theorem}{Theorem}[section]
\newtheorem{lemma}[theorem]{Lemma}
\newtheorem{proposition}[theorem]{Proposition}
\title{Copied retained centers in two-stage multiplication networks\\
\large Exact corner blocks and a conditional fixed-tape bound}
\author{icekylinx\\\small With substantial OpenAI GPT-6 Astra and Codex assistance}
\date{8 October 2026}
\begin{document}\maketitle
\begin{abstract}
A retained center can supply its scatter from a transformed copy while
its original role proceeds directly to cleanup. Paying the copy transform
reduces its local rank charge from $r+h$ to $r+(h-r)$ in both invocation
orientations. Combining this schedule with two tensor stages, a certified
$(25,23)$ bit corner and a $(28,28)$ mixed-center complex network gives
$T(n)=O(n(\log n)^{1-\kappa})$, where
$\kappa=384569/10^{10}=3.84569\times10^{-5}$.
The result is conditional on the retained analytic and fixed-tape
interfaces. Exact finite certificates support counts and inequalities;
the all-size conclusion also uses the written mathematical and tape proofs.
\end{abstract}

\section{Scope and attribution}
The model is a fixed finite-alphabet multitape Turing machine. Every
finite producer, basis, address prime, table and tape count is fixed
independently of input length. The bit scalar payload is $\mathbb F_2$;
its rational address geometry is later reduced at an eligible odd prime.
The complex payload is Gaussian dyadic with binary address geometry.
These distinct fields and their physical role budgets are kept separate.
Constructive setup and the analytic, exact-recovery and prime-existence
thresholds remain eventual assumptions of the retained framework.

The present copied-center schedule, its complex transfer, the selected
finite specialization and circuit-specific assembly are contributed by
icekylinx with substantial OpenAI GPT-6 Astra assistance; integration
uses Codex assistance. The immediate predecessor is icekylinx's PR32,
\texttt{0ef3aeb61f55cc0b321ce6a0ef00acee25cefe52}, with proof
\path{notes/structured-bulk-note.tex}; earlier own contributions are
PR24 endpoint gauges, PR18 partial swaps, and PR10 projector batching.
Their retained producers, labels, mixed centers, whole-residual complex
normal form, integer-width rows and semantic/bulk composition are used
through their proved interfaces, with only the changes specified here.

Aurel Prosz (Paureel), at
\texttt{c82d09eb4781b68435e3fa3eb6beea72fa900fab}, supplies the two-stage
topology and paid endpoint-copy correction \cite{paureel}. Swapnil Jain's
linked two-stage development, commit
\texttt{ae405eb474d1486b2d8aef90139869f927f4836a}, is credited as in
PR29. Zhihao Chen's PR29 \cite{pr29}, pinned at
\texttt{9d963275075fa98f1da821e27757b238dafd6b3c}, supplies the unequal
two-stage compatibility, common basis, gauges and physical accounting.
Rohan Arun's PR31 \cite{pr31}, method commit
\texttt{02f68f95369c39955c5ffcfddbd75cd83e26b305}, supplies the
incidence factorization, rank-partition zero identities and rational
pivot-witness method. Dominik Scholz's PR33 \cite{pr33}, pinned at
\texttt{a499a345c040aa15b68e984a0904eadc485c7a7e}, supplies its portable
dimension parameterization, here specialized to $(25,23)$.

The semantic child method, arbitrary-source routing, phase-cell inverse
and bulk locality/tape arguments are RaD/hipotures, through PR20 at
\texttt{45d9b60355872041f9275f77c057514def0d45bc}; Zhihao Chen's PR23
supplies the circuit-specific composition, and PR21 the translated
complex auxiliary endpoints. Their proof map and pinned files remain in
\path{notes/structured-bulk-assembly.tex} and
\path{references/semantic-bulk/}. The Gaussian error bounds, including
computable downward approximations and logarithmic inverse normalization,
remain those of that retained analytic transfer.

Douglas Colkitt's framework, the upstream OpenAI manuscript, eumemic's
source frames and complex/resampling work, Bortlesboat's alignment,
dleen's retained totals, the earlier finite work of Zhihao Chen and
Rohan Arun, and Harvey--van der Hoeven's analytic background retain their
credits in \texttt{NOTICE} and \texttt{SOURCES.json}. PR29/31's recorded
OpenAI Codex assistance, PR33's recorded Anthropic Claude Opus 5.5 and
OpenAI GPT-6 Astra/Codex assistance, and the RaD source's recorded AI
assistance remain attributed to their sources. Imported licenses and
notices are preserved under \path{references/copied-centers/}.

\section{Copied retained-center reads}
\input{notes/copied-centers-lemma.tex}
\section{Selected bit network and exact common corner}
\input{notes/copied-centers-bit.tex}
\section{Selected complex network and semantic assembly}
\input{notes/copied-centers-complex.tex}
\input{notes/copied-centers-assembly.tex}

\begin{thebibliography}{9}
\bibitem{paureel} Aurel Prosz (Paureel). Two-stage topology and charged
endpoint correction. \url{https://github.com/Paureel/integer-mult-bounds}.
Retained proof: \path{references/copied-centers/pr29/paureel-direct-transfer.tex}.
\bibitem{pr29} Zhihao Chen (jacklightChen). Unequal two-stage construction.
\url{https://github.com/CrocSwap/integer-mult-bounds/pull/29}.
Proof: \path{references/copied-centers/pr29/two-stage-16-note.tex}.
\bibitem{pr31} Rohan Arun (rohanarun). Contiguous two-stage data corners.
\url{https://github.com/CrocSwap/integer-mult-bounds/pull/31}.
Method provenance: \path{references/copied-centers/pr33/SOURCE.json}.
\bibitem{pr33} Dominik Scholz. Parameterized two-stage corner specialization.
\url{https://github.com/CrocSwap/integer-mult-bounds/pull/33}.
Proof: \path{references/copied-centers/pr33/two-stage-corners-47-45-note.tex}.
\bibitem{pr32} icekylinx. Structured blocks and semantic bulk assembly.
\url{https://github.com/CrocSwap/integer-mult-bounds/pull/32}.
\bibitem{pr23} Zhihao Chen. Semantic precision and bulk composition.
\url{https://github.com/CrocSwap/integer-mult-bounds/pull/23}.
\bibitem{rad} RaD project (hipotures). Analytic and fixed-tape constructions.
\url{https://github.com/hipotures/rad/tree/45d9b60355872041f9275f77c057514def0d45bc/research/integer-multiplication-bounds}.
\end{thebibliography}
\end{document}
